% Benchmark Methodology
\section{Benchmark Methodology}
\label{sec:methodology}

\subsection{Fairness}
All engines execute identical logical queries via a common
adapter interface.
Query translation overhead is excluded from runtime measurements.
Each engine is configured with default settings;
no engine-specific tuning is performed.

\subsection{Cache Modes}
\textbf{Cold cache:} The engine state is re-initialized before
each run. Ten runs are collected.
\textbf{Warm cache:} The engine retains state across runs.
Fifty runs are collected after three warm-up iterations.

\subsection{Outlier Handling}
Outliers are discarded only when Grubbs' test identifies them
as statistically significant ($\alpha = 0.05$).
All discarded outliers are documented.

\subsection{Correctness Verification}
Before recording performance, every benchmark run verifies
cross-engine result equivalence.
If any engine produces a different canonical result for the
same query, the benchmark is rejected and flagged.

\subsection{Statistical Analysis}
We report mean, median, standard deviation,
P50, P95, P99 percentiles,
95\% confidence intervals using the normal approximation,
and bootstrap confidence intervals (10,000 resamples).

\subsection{Reproducibility}
Every benchmark run records:
Python version, uv lock hash, Git commit hash,
CPU model, memory, operating system,
random seed, dataset checksum, engine versions,
and timestamp.
All generated figures and tables include the
reproducibility hash in their caption.
